\documentclass[12pt]{article}
\usepackage[utf8]{inputenc}
\usepackage[T1]{fontenc}
\usepackage{hyperref}
\title{Central Tagbanwa Language Resources for K-12 Curriculum Development}
\author{Compiled by Claude Sonnet 5 (claude-sonnet-5) for Andr\'as Kornai}
\date{}
\begin{document}
\maketitle

\section{Introduction}
Central Tagbanwa (ISO 639-3 tgt; Glottocode cent2090) is a Palawanic Austronesian language spoken in northwestern Palawan, Philippines, in Taytay, Roxas and San Vicente. The UP Linguistics capsule calls it nearly extinct (EGIDS 8b) and reports about 2,000 users in 1985.

\section{Primary Educational Resources}
No primers, textbooks or school programmes were found.

\section{Linguistic and Reference Materials}
\begin{itemize}
\item UP Linguistics language capsule for Central Tagbanwa: \url{https://linguistics.upd.edu.ph/the-katig-collective/language-capsules/tagbanwa-central/}
\item Wikipedia, with phonology and pronoun tables: \url{https://en.wikipedia.org/wiki/Central_Tagbanwa_language}
\item Glottolog entry: \url{https://glottolog.org/resource/languoid/id/cent2090.bigmap.html}
\end{itemize}
The search summary names R. Scebold (2003), \emph{Central Tagbanwa: A Philippine Language on the Brink of Extinction} (Linguistic Society of the Philippines), as the main scholarly work, and says SIL hosts a PDF. I did not see a URL for it, so it is not listed as a verified resource.

\section{Religious and Community Resources}
None were found.

\section{Digital and Audio Resources}
\begin{itemize}
\item Omniglot page on Central Tagbanwa, covering the Latin script and the Tagbanwa script: \url{https://www.omniglot.com/writing/centraltagbanwa.htm}
\item Endangered Languages Project entry: \url{https://www.endangeredlanguages.com/lang/5087}
\end{itemize}
No audio specific to Central Tagbanwa was found. The Global Recordings result concerned a different variety.

\section{Curriculum Development Context}
The Tagbanwa script and the Scebold monograph give a starting point. With speakers mostly elderly, a revitalization curriculum would depend on documenting those speakers soon.
\end{document}